\subsection{Symmetry effects by graph family (RQ4)}
\label{sec:symmetry-summary}
The archived $L(2,1)$ experiment compares a baseline with a symmetry
configuration on $C_3,\ldots,C_{50}$ and seven product families with
factor sizes $3\leq n,m\leq10$. It uses Glucose, a fresh solver at each
span trial, and limits of 60 seconds for cycles and 180 seconds for
products. Each graph/configuration is measured once; execution order
alternates. Both configurations apply the same encoder preprocessing.
These timings are separate from the 30-second CaDiCaL--Gurobi experiment.

\begin{table}[htbp]
\centering\small
\caption{Symmetry effects in the paired $L(2,1)$ experiment.}
\label{tab:compact-symmetry}
\input{\SubmissionSummaryDir/symmetry.tex}
\tablenote{B and S denote baseline and symmetry. R fixes the cycle root
to 0; O orders a reflected pair;
-- adds no family-specific constraint. OPT/pairs counts jointly optimal
pairs out of all tested pairs. On these pairs, $\Delta V$ and $\Delta C$
are median percentage reductions in allocated variables and emitted
clauses, measured after common encoder preprocessing at the same span.
The next column counts fewer/equal/more clauses. The last column is the
median of paired time ratios; values above 1 favor symmetry. Unresolved
pairs remain in the denominator but are excluded from the medians.}
\end{table}

Both configurations report the same optimal span on 494 of 496 pairs.
The two unresolved pairs are $C_9\square C_{10}$ and
$C_{10}\square C_9$, with intervals $[5,10]$ and $[7,8]$, respectively,
under both configurations. Coordinate exchange is an isomorphism, so
these are two presentations of the same graph. Their different retained
incumbents are not different optima; vertex order and search history
can affect the search. All four unresolved runs have incomplete search
statistics and are excluded from the descriptive timing and counter
summaries. For cycles, root fixing with reflection ordering
reduces the median clause count by 12.58\%, but the median time ratio is
1.000 to three decimals. In the four product families using an order
constraint, median clause reductions range from 0.58\% to 1.14\% and
median time ratios from 1.021 to 1.120. The three families with no added
constraint have identical model counts; their timing differences describe
repeat variation. Thus the data show family-dependent reductions in CNF
size, without a proportional runtime reduction. Single measurements and
the combined cycle constraints do not support significance claims or
separate attribution of the cycle effect to root fixing and ordering.

\begin{table}[htbp]
\centering\small
\caption{Search counters in complete symmetry-run pairs.}
\label{tab:compact-counters}
\input{\SubmissionSummaryDir/symmetry_counters.tex}
\tablenote{B and S denote baseline and symmetry. Each entry sums the
counters over all completed span trials and all complete jointly OPT
runs in that family; these are totals, not medians. Different runs can
visit different spans, so this is not a fixed-span ablation. Partial
counters, including zeros after an interrupted trial, are excluded.}
\end{table}

Across the cycle runs, conflicts fall from 490 to 16 and decisions from
1014 to 304, while the median time ratio remains near one.
Aggregate conflicts fall by about 45\% for $C_n\square C_m$ and 32\%
for $C_n\circ C_m$. Counters describe completed solver search;
wall time also includes encoding, startup, validation and cleanup.
Reduced counts can therefore coexist with little change in total time.
The measurements do not isolate which overhead explains the difference
or establish a causal relationship between clause count and runtime.
